% SECTION: Principal Component Analysis (PCA)
\section{Principal Component Analysis (PCA)}

% SOLUTION
\subsection*{Lösung zu Aufgabe \ref{task:pca-derivation-eigenvalue-problem}}

Wir stellen zunächst die Lagrange-Funktion auf. Diese lautet
\begin{equation*}
	\lagrange(\bu_1, \lambda_1)
		= \bu_1^\top \bS \bu_1 + \lambda_1 \bigl( \bu_1^\top \bu_1 - 1 \bigr).
\end{equation*}
Die partielle Ableitung von $\lagrange$ nach $\bu_1$ ist
\begin{equation*}
	\frac{\partial}{\partial \bu_1} \lagrange(\bu_1, \lambda_1)
		= 2 \bS \bu_1 - 2 \lambda_1 \bu_1.
\end{equation*}
Wir setzen Null und lösen nach $\bu_1$ auf. Wir erhalten schließlich das Eigenwertproblem
\begin{equation*}
	\bS \bu_1 = \lambda_1 \bu_1.
\end{equation*}


% SOLUTION
\subsection*{Lösung zu Aufgabe \ref{task:pca-algorithm}}


% SOLUTION
\subsection*{Lösung zu Aufgabe \ref{task:pca-minimum-error}}




\newpage